\section{Related Work}\label{sec:related work}

Error-bounded lossy compression reduces scientific data volumes while providing explicit control over reconstruction error~\cite{zhao2021optimizing,di2025survey}. In scientific data-management workflows, it can lower storage and I/O costs and accelerate wide-area data transfer~\cite{cappello2019usecases,liu2023ocelot,shah2026solver}. The magnitude of these practical benefits depends in part on the achieved data reduction, which in turn depends on how well the compression model matches the organization of the input data. We therefore review existing methods for structured arrays, unstructured data, and high-order element data.

\subsection{Structured Tensor Compression}

Many error-bounded compressors target scientific data stored as regular multidimensional arrays~\cite{di2025survey}. Their regular layout makes neighboring values inexpensive to locate and enables efficient prediction or transform schemes. The SZ family follows a prediction--quantization--encoding pipeline. SZ2~\cite{tao2017significantly} uses predictors such as Lorenzo prediction and blockwise linear regression, whereas SZ3~\cite{zhao2021optimizing,liang2023sz3} uses multilevel interpolation along each dimension of an array. QoZ~\cite{liu2024high} builds on this interpolation approach by exploring a broader set of prediction configurations and using autotuning to select among them. Other compressors use transforms. ZFP~\cite{lindstrom2014fixed} applies a decorrelating transform to small fixed-size blocks, and SPERR~\cite{li2023lossy} applies a wavelet transform before encoding the transformed data.

These compressors can process high-order element data after simple preprocessing, but their models see only the resulting array layout. They do not account for the nonuniform point spacing within an element or the statistical differences among element-local positions.

\subsection{Unstructured Data Compression}

Recent error-bounded compressors for unstructured scientific data construct predictors from relationships in physical space. Ren et al.~\cite{ren2024prediction} propose a prediction--traversal approach that arranges mesh nodes into sequences and predicts newly visited nodes from their neighbors. UMC~\cite{wu2025enabling} preprocesses mesh information and provides several predictors for different unstructured meshes. Gong et al.~\cite{gong2024general} introduce a multi-component framework that interpolates mesh data onto a regular grid and separately compresses the interpolated data and residuals.

Such methods are designed for data whose useful relationships must be recovered from coordinates or mesh connectivity. High-order element data also have irregular physical geometry, but the points within every element follow the same quasi-structured arrangement in a reference space. Moreover, corresponding local positions have similar statistical characteristics across elements, even though these characteristics may differ from one local position to another. This cross-element regularity need not appear as continuity in the stored element order and can therefore be difficult to identify directly. General-purpose unstructured compressors do not explicitly use this combination of intra-element geometry and cross-element regularity.

\subsection{Compression for High-Order Element Data}\label{sec:related work:prior work}

Several studies have considered compression directly in high-order element workflows. Shah et al.~\cite{shah2026solver} integrate SZ3 into the NekRS parallel output workflow. Before writing an output, each parallel process compresses the simulation fields stored in its local spectral elements, thereby reducing the amount of data sent to storage. They also analyze how errors introduced by compression affect values and derivatives computed by the spectral-element solver. However, they use SZ3's existing prediction algorithm rather than adapting its compression model to nonuniform GLL spacing or the statistical differences among element-local positions.

Marin et al.~\cite{marin2016spectral} develop a specialized compressor for Nek5000 data defined on Gauss--Lobatto--Legendre or Chebyshev nodes. It uses the basis-matched transform---the discrete Legendre transform (DLT) for GLL nodes and the discrete Chebyshev transform (DChT) for Chebyshev nodes---to represent each element in coefficient space, sets selected coefficients to zero, and encodes the resulting zeros compactly. However, the decompressed data are not guaranteed to satisfy a pointwise error bound. Nevertheless, its observation that element-wise transforms concentrate values near zero motivates the coefficient representation used by MRCZ.
